\documentclass[a4paper]{article}
\usepackage[pdftex]{graphicx}
\usepackage[utf8]{inputenc}
\usepackage{enumerate}
\usepackage{amssymb}
\usepackage{eurosym}
\usepackage{tikz}
\usepackage{href-ul}
\hypersetup{
	colorlinks=true,
	linkcolor=blue,
	urlcolor=blue}
\usepackage{geometry}
\geometry{a4paper, top=15mm, left=15mm, right=15mm, bottom=15mm,
	headsep=10mm, footskip=12mm}

\begin{document}
	\begin{tikzpicture}(10,3)
		\draw[ultra thick](2,0) --(10,0) -- (10,1.3) --(2,1.3) -- (2,0);
		\draw[fill=black](2,0)-- (0,.6) -- (2,.6) -- (2,0);
		\node at (6,.6) {\large Solution sheet from \href{http://www.okuyakl.com}{www.okuyakl.com}};
	\end{tikzpicture}
	\vspace{0.5 cm}
	
	\noindent{\bf Task 1. a)}\\
	For $E_1$="Sum  = 2" there are the following options:
	$$E_1=\{20;02;11\}$$
	The probability of a die showing a 2 is $1 \over 6$, for a 1 it is $1 \over 2$, and for a 0 it is $1 \over 3$. So the probability for the sum of the points is 2:
	$$p(E_1)= 2 \cdot {1 \over 6} \cdot {1 \over 3} + \left({1 \over 2}\right)^2={13\over 36}$$
	
	\noindent{\bf Task 1. b)}\\
	$$
	\renewcommand{\arraystretch}{2}
	\begin{array}{c|c|c|c|c|c}
		x_i & 0 & 1 & 2 & 3 & 4 \\
		\hline
		P(X=x_i) & {1\over9 } & {1\over 3} & \frac{13}{36} & {1\over 6} & {1\over 36}\\
	\end{array}
	$$
	
	\noindent{\bf Task 1. c)}\\
	For $E_2$="3 balls of the same color" there are the following options:
	$$E_2=\{rrr;ggg;bbb\}$$
	The probability of this is:
	$$p(E_2)= 3 \cdot {3 \over 9} \cdot {2\over8} \cdot {1\over 7}= {1\over 28}$$
	
	\noindent{\bf Task 1. d)}\\
	For this we create a term for the expected value from the perspective of the transmission line:
	$$
	\renewcommand{\arraystretch}{2}
	\begin{array}{rcll}
		E(x)&=& (2-B) \cdot p(E_2)+ 2 \cdot (1-p(E_2))\\
		1.25 &=& {2\over 28}- B\cdot {1\over 28} +{54\over 28}\\
		B &=& 0.75 \cdot 28 = 21\\
	\end{array}
	$$
	The profit must be \EUR{21}.
	
	\noindent{\bf Task 2. a)}\\
	
	The probability for event (i) corresponds to a Bernoulli chain of length 10:
	
	$$P_{1 \over 3}^{10}(X=5)=B(10;{1 \over 3}; 5)={10 \choose 5} \cdot \left( {1 \over 3} \right)^5 \cdot
	\left( {2 \over 3} \right)^5=13.6\%$$
	
	For the probability of event (ii) we transform and read from the table:
	
	$$ P_{1 \over 3}^{10}(4\le X \le 8)= F_{1 \over 3}^{10}(8)-F_{1 \over 3}^{10}( 3)=0.999644-0.559264=44.0\%$$
	
	\noindent {\bf Task 2. b)}
	
	We create the following table of probability distribution:
	
	$$
	\begin{array}{c|c|c|c}
		X_i & -2.90 & 0.10 & 1.10 \\
		\hline
		P(X_i) & {1 \over 3} & {1 \over 6} & {1 \over 2}\\
	\end{array}
	$$
	
	The expected value $E(x)$ of the random variable “profit” is:
	
	$$E(x)= -2.90 \cdot {1 \over 3} + 0.1 \cdot {1 \over 6} + 1.1 \cdot {1 \over 2}=-0.4$$
	
	So you make an average loss of \EUR{0.4} per game.
	
	\noindent {\bf Task 3. a)}
	
	The probability we are looking for is again a Bernoulli chain:
	
	$$P_{0.03}^{12}(X=2)=B(12,0.03, 2)={12 \choose 2} \cdot 0.03^2 \cdot 0.97^{10 }=4.4\%$$
	
	\newpage
	\noindent {\bf Task 3. b)}
	
	This is a so-called three times minimum task: We work with the counter event:
	
	$$
	\begin{array}{rcl}
		1-P_{p}^{12}(X=0) & \ge & 0.9\\
		1-(1-p)^{12}& \ge & 0.9\\
		(1-p)^{12} & \le & 0.1 \quad | \sqrt[12]{}\\
		1-p & \le & 0.8254 \\
		p & \ge & 0.18 \quad \textnormal{also} \quad p \ge 18\%
	\end{array}
	$$
	
	\noindent {\bf Task 4.}
	The probability distribution for $Y$
	is:
	$$
	\renewcommand{\arraystretch}{2}
	\begin{array}{|c|c|c|c|}
		\hline
		y_i & 1450 & 1900 & 2350 \\
		\hline
		P(Y=y_i)& 0.6 & 0.1 & 0.3 \\
		\hline
	\end{array}
	$$
	The expected value here is:
	$$E(Y)=1450\cdot 0.6 + 1900 \cdot 0.1 +2350 \cdot 0.3 = 1765$$
	The variance is:
	$$Var(Y)=(1450-1765)^2\cdot 0.6 +(1900-1765)^2\cdot 0.1+(2350-1765)^2\cdot 0.3 =164025$$
	So the standard deviation is:
	$$\sigma(Y)=\sqrt{164025}=405$$
	
	\begin{center}
		\includegraphics[width=7 cm]{../../viecher/eendcomic.pdf}
		
		Here you can go back to the \href{https://www.okuyakl.de/math/m12stoL064/ae064.pdf}{task sheet}
	\end{center}
	
\end{document}